\chapter{Probability and Randomized Algorithms}

\providecommand{\E}{\ensuremath{\mathbb{E}}}

So far, every algorithm you have met has been a predictable little machine: feed it
the same input twice and it does exactly the same thing twice, step for step. In
this chapter we hand the machine a coin. We let it flip that coin and decide what
to do based on the result. That sounds reckless --- who wants a sorting routine
that behaves differently each time you run it? --- but it is one of the most
useful ideas in algorithmic design. Randomness, used
deliberately, can make an algorithm \emph{faster on average}, \emph{simpler to
reason about}, and \emph{immune to the adversary who picks your worst-case
input}.

To understand why, we need two things. First, a working vocabulary of
probability --- not a full course on probability, just the handful of tools that algorithm
designers reach for constantly. Second, a way to talk about ``how long does this
take \emph{on average}'' that is just as rigorous as the worst-case analysis you
already know. Both threads come together in a single
device called the \emph{indicator random variable}. Master that one trick and a
huge swath of average-case analysis becomes almost mechanical.

\begin{focusbox}
The destination of this chapter: you will be able to take an algorithm whose
running time depends on chance (or on the order of its input), set up the
expected cost with indicator random variables, and \emph{compute} that
expectation rather than guess at it. We build the probability tools first, then
spend them on three concrete problems: the hiring problem, randomly permuting an
array, and the famous birthday-paradox style counting that underlies hashing and
quicksort.
\end{focusbox}

\section{Two Different Kinds of ``Average''}

Before any formulas, let us get the big picture straight, because beginners
routinely confuse two genuinely different ideas that both get called
``average-case.''

\begin{intuitionbox}
There are two distinct sources of uncertainty an analysis might average over.

\textbf{(1) Random input.} The algorithm is completely deterministic, but we
\emph{assume} the input arrives drawn from some probability distribution (say,
``all $n!$ orderings of the array are equally likely''). We average the running
time over that distribution. This is \emph{probabilistic analysis} of a
deterministic algorithm.

\textbf{(2) Random choices inside the algorithm.} The input is whatever it is ---
possibly chosen by a malicious adversary --- but the \emph{algorithm itself}
flips coins as it runs. We average over the algorithm's own coin flips. This is a
\emph{randomized algorithm}.

The second is far more satisfying, because \emph{you} control the coins. You do
not have to cross your fingers and hope the input is nicely shuffled; you shuffle
it yourself. The guarantee then holds for \emph{every} input.
\end{intuitionbox}

The hiring problem below starts as case (1) and we will deliberately convert it
into case (2). That conversion --- turning an assumption about the input into a
fact you enforce with your own randomness --- is the single most important habit
in this chapter.

\section{The Hiring Problem}
\label{ch05:hiring}

Here is a story with a clean cost structure. You are hiring an office assistant
through an employment agency. Each day the agency sends you exactly one
candidate. Interviewing a candidate is cheap (cost $c_i$ per interview), but
\emph{hiring} is expensive: to hire someone you must first fire whoever currently
holds the job and pay a large fee $c_h$ for the new hire ($c_h \gg c_i$).

You want the best possible assistant, so you adopt a simple rule:

\begin{algorithm}[H]
\caption{\textsc{Hire-Assistant}$(n)$}
\KwIn{a sequence of $n$ candidates, interviewed one at a time}
\KwOut{the running staff is always the best candidate seen so far}
$best \gets 0$ \tcp*{candidate $0$ is a fictional least-qualified dummy}
\For{$i \gets 1$ \KwTo $n$}{
  interview candidate $i$\;
  \If{candidate $i$ is better than candidate $best$}{
    $best \gets i$\;
    hire candidate $i$\;
  }
}
\end{algorithm}

You interview everyone: that costs $n\,c_i$, no surprises, it happens every time.
The interesting cost is hiring. Every time you meet someone better than everyone
before, you hire them (and eat the fee $c_h$). \emph{How many times do you hire?}

\begin{intuitionbox}
The worst case is easy and brutal. If the candidates arrive in strictly
increasing order of quality --- worst first, best last --- then \emph{every}
candidate is better than all the predecessors, so you hire all $n$ of them. Total
hiring cost $n\,c_h$. Ouch.

But surely that exact adversarial order is rare. If candidates show up in a
\emph{random} order, you would expect to hire only occasionally: most new
candidates are not record-breakers. The central question is: how many hires occur, on
average?
\end{intuitionbox}

The answer, which we will prove, is shockingly small: about $\ln n$ hires. For
$n = 1000$ candidates you expect to hire roughly $7$ times, not $1000$. To get
there we need the right tool.

\section{Random Variables and Expectation}

Let us nail down the vocabulary. A \emph{sample space} $S$ is the set of all
possible outcomes of a random experiment, and a \emph{probability distribution}
assigns to each outcome a number in $[0,1]$ so that the numbers sum to $1$. An
\emph{event} is a subset of $S$, and its probability is the sum of the
probabilities of the outcomes inside it.

\begin{definition}[Random variable]
\label{ch05:def-rv}
A \emph{random variable} $X$ is a function from the sample space to the real
numbers. It attaches a number to each outcome. We write $\Pr\{X = x\}$ for the
probability of the event ``$X$ takes the value $x$,'' meaning the total
probability of all outcomes mapped to $x$.
\end{definition}

The number we care about most is the \emph{average value} of a random variable,
weighting each possible value by how likely it is.

\begin{definition}[Expectation]
\label{ch05:def-exp}
The \emph{expected value} (or \emph{expectation}, or \emph{mean}) of a random
variable $X$ is
\[
  \E[X] \;=\; \sum_{x} x \cdot \Pr\{X = x\},
\]
the sum over all values $x$ that $X$ can take, of the value times its
probability.
\end{definition}

\begin{example}
\label{ch05:ex-die}
Roll a fair six-sided die; let $X$ be the number shown. Each value has
probability $\tfrac16$, so
\[
  \E[X] = 1\cdot\tfrac16 + 2\cdot\tfrac16 + 3\cdot\tfrac16
        + 4\cdot\tfrac16 + 5\cdot\tfrac16 + 6\cdot\tfrac16
        = \frac{21}{6} = 3.5.
\]
Notice that $3.5$ is not a face of the die at all. The expectation is a
weighted center of mass, not a value you ever literally observe. That is normal
and important: expectations need not be ``achievable'' outcomes.
\end{example}

The single property of expectation that powers nearly everything in this chapter
is its \emph{linearity}.

\begin{theorem}[Linearity of expectation]
\label{ch05:thm-linearity}
For any two random variables $X$ and $Y$ defined on the same sample space, and
any constant $a$,
\[
  \E[X + Y] = \E[X] + \E[Y],
  \qquad
  \E[aX] = a\,\E[X].
\]
More generally, for any random variables $X_1, \dots, X_n$,
\[
  \E\!\left[\sum_{i=1}^{n} X_i\right] = \sum_{i=1}^{n} \E[X_i].
\]
\end{theorem}

\begin{focusbox}
Read the last line again and notice the missing fine print. \textbf{There is no
independence assumption.} Linearity of expectation holds even when the $X_i$
push and pull on each other in complicated ways. This is the loophole that makes
hard problems easy: we get to break a tangled quantity into a sum of simple
pieces, take the expectation of each piece separately, and add up --- without
ever worrying whether the pieces are correlated.
\end{focusbox}

\begin{trapbox}
Linearity is for \emph{sums}, not products. In general
$\E[XY] \neq \E[X]\,\E[Y]$; that product rule needs $X$ and $Y$ to be
\emph{independent}. Beginners who get drunk on linearity often try to factor a
product of random variables and quietly produce nonsense. If you find yourself
splitting a product, stop and check independence explicitly.
\end{trapbox}

\section{Indicator Random Variables: The Master Trick}
\label{ch05:indicators}

Here is the device that simplifies average-case calculations. Suppose you want to count
how many times some event happens --- how many record-breaking candidates, how
many collisions, how many inversions. Counting is hard to do directly because the
events interact. The trick is to attach a tiny $0/1$ variable to \emph{each}
possible occurrence, then add them up.

\begin{definition}[Indicator random variable]
\label{ch05:def-indicator}
Given an event $A$, the \emph{indicator random variable} $I\{A\}$ is defined by
\[
  I\{A\} =
  \begin{cases}
    1 & \text{if event } A \text{ occurs},\\
    0 & \text{if event } A \text{ does not occur}.
  \end{cases}
\]
\end{definition}

The bridge between probability and counting is captured in a single concise lemma:

\begin{lemma}[Expectation of an indicator]
\label{ch05:lem-indicator}
For any event $A$, the indicator $X_A = I\{A\}$ satisfies
\[
  \E[X_A] = \Pr\{A\}.
\]
\end{lemma}

\begin{proof}
$X_A$ takes only the values $1$ and $0$. By the definition of expectation,
\[
  \E[X_A] = 1\cdot\Pr\{A\} + 0\cdot\Pr\{\overline{A}\} = \Pr\{A\}.
\qedhere
\]
\end{proof}

\begin{intuitionbox}
The expected value of a yes/no variable is just the probability that the answer
is yes. So an indicator is a bridge: it converts a \emph{probability} (something
you compute by reasoning about chance) into an \emph{expectation} (something you
can add up with linearity). The recipe, which we will run three times in this
chapter, is:

\begin{enumerate}
  \item Write the quantity you want to count as a \textbf{sum of indicators},
        one per possible occurrence.
  \item Take the expectation, and use \textbf{linearity} to push it inside the
        sum.
  \item Replace each $\E[\text{indicator}]$ by the corresponding
        \textbf{probability} (Lemma~\ref{ch05:lem-indicator}).
  \item Compute those probabilities --- usually each one is easy --- and
        \textbf{sum}.
\end{enumerate}
\end{intuitionbox}

Figure~\ref{ch05:fig-recipe} summarizes this pipeline: a hard count on
the left, dissolved into easy probabilities on the right.

\begin{figure}[H]
\centering
\begin{tikzpicture}[
    node distance=6mm,
    box/.style={draw,rounded corners,align=center,inner sep=4pt,minimum height=1cm},
    arr/.style={-Stealth,thick},
  ]
  \node[box,fill=trapred!10]    (count) {count\\$X=\sum_i X_i$};
  \node[box,fill=sourceblue!10,right=of count] (lin) {linearity\\$\E[X]=\sum_i\E[X_i]$};
  \node[box,fill=quotientorange!12,right=of lin] (ind) {indicator\\$\E[X_i]=\Pr\{A_i\}$};
  \node[box,fill=targetgreen!14,right=of ind] (sum) {sum the\\probabilities};
  \draw[arr] (count) -- (lin);
  \draw[arr] (lin) -- (ind);
  \draw[arr] (ind) -- (sum);
\end{tikzpicture}
\caption{The indicator recipe as a pipeline. Each arrow is one of the four steps;
the only place that needs real thought is the rightmost box, computing each
$\Pr\{A_i\}$, and those are usually one-liners.}
\label{ch05:fig-recipe}
\end{figure}

\subsection*{Warm-up: how many heads in $n$ coin flips?}

Flip a fair coin $n$ times; let $X$ be the total number of heads. You could grind
through the binomial distribution, but watch how the indicator recipe dissolves
it. Let $X_i = I\{\text{flip } i \text{ is heads}\}$. Then $X = \sum_{i=1}^n X_i$,
and
\[
  \E[X] = \E\!\left[\sum_{i=1}^n X_i\right]
        = \sum_{i=1}^n \E[X_i]
        = \sum_{i=1}^n \Pr\{\text{flip } i \text{ heads}\}
        = \sum_{i=1}^n \tfrac12
        = \frac{n}{2}.
\]
No binomial coefficients, no summation identities. Four lines. That is the power
you are buying.

\begin{labbox}
Trace the recipe by hand for $n = 3$ flips. List the $8$ equally likely outcomes
(HHH, HHT, \dots, TTT), count the heads in each, average the eight counts
directly, and confirm you get $3/2 = 1.5$. Then redo it with the indicator
method and feel how much shorter the second path is. Once both give $1.5$ you
have internalized the trick.
\end{labbox}

\section{Solving the Hiring Problem with Indicators}
\label{ch05:hiring-solved}

Now we cash in. Assume the candidates arrive in a \emph{uniformly random order}:
all $n!$ orderings are equally likely (this is the probabilistic-analysis
assumption; we will justify or replace it soon). Let $X$ be the number of times
we hire. We want $\E[X]$.

\textbf{Step 1: write $X$ as a sum of indicators.} Define, for each candidate
$i$,
\[
  X_i = I\{\text{candidate } i \text{ is hired}\}.
\]
Hiring happens at step $i$ exactly when candidate $i$ is better than all of
candidates $1, 2, \dots, i-1$. Then $X = X_1 + X_2 + \dots + X_n$.

\textbf{Step 2: linearity.}
\[
  \E[X] = \sum_{i=1}^n \E[X_i] = \sum_{i=1}^n \Pr\{\text{candidate } i
  \text{ is the best of the first } i\}.
\]

\textbf{Step 3: compute the probability.} Look only at the first $i$ candidates.
Because the order is uniformly random, each of these $i$ candidates is equally
likely to be the best \emph{among those first $i$}. There are $i$ of them, so the
chance that the best one happens to sit in position $i$ is exactly $1/i$. Hence
$\E[X_i] = 1/i$.

\begin{trapbox}
Notice we only needed the \emph{relative} order of the first $i$ candidates to be
uniform, and by symmetry every candidate is equally likely to be the local
maximum. A common error is to instead reason ``candidate $i$ must be the global
best, probability $1/n$.'' That is wrong --- we are not asking whether $i$ is the
best overall, only whether $i$ beats the $i-1$ people who came before. Those are
different events with different probabilities.
\end{trapbox}

\textbf{Step 4: sum.}
\[
  \E[X] = \sum_{i=1}^n \frac{1}{i} = H_n,
\]
the $n$-th \emph{harmonic number}. And we know from calculus (the harmonic series
is squeezed between two logarithms) that
\[
  \ln n \;<\; H_n \;\le\; \ln n + 1,
  \qquad\text{so}\qquad
  \E[X] = \ln n + O(1).
\]

\begin{formalbox}
\textbf{Theorem (expected hires).} Assuming candidates are interviewed in a
uniformly random order, \textsc{Hire-Assistant} hires an expected $H_n = \ln n +
O(1)$ times, for a total expected hiring cost of $O(c_h \ln n)$.
\end{formalbox}

Let us make that concrete.

\begin{example}
\label{ch05:ex-hiring-numbers}
For $n = 1000$: $H_{1000} \approx \ln 1000 + 0.577 \approx 6.91 + 0.58 \approx
7.49$. So you expect to hire about $7$ or $8$ times, not $1000$. For $n = 10^6$:
$H_{10^6} \approx \ln 10^6 \approx 13.8$. A million applicants, fourteen hires.
The worst case wanted a million hires; the average wants fourteen. That gap is
the entire reason this chapter exists.
\end{example}

\subsection*{Randomizing the algorithm}

There is a catch hiding in the phrase ``assuming a uniformly random order.''
Real candidates do not arrive in random order; the agency might, out of malice
or laziness, send them worst-first and force $n$ hires every single time. The
probabilistic analysis describes a hope, not a guarantee.

So we \emph{enforce} the hope. Before interviewing anyone, we collect all the
applications and \emph{randomly permute the list ourselves}, then interview in
that scrambled order:

\begin{algorithm}[H]
\caption{\textsc{Randomized-Hire-Assistant}$(n)$}
\KwIn{a list of $n$ candidates in arbitrary (possibly adversarial) order}
\KwOut{expected number of hires is $H_n$, for \emph{every} input}
randomly permute the list of candidates\;
$best \gets 0$\;
\For{$i \gets 1$ \KwTo $n$}{
  interview candidate $i$\;
  \If{candidate $i$ is better than candidate $best$}{
    $best \gets i$; \quad hire candidate $i$\;
  }
}
\end{algorithm}

Now the random order is a \emph{fact we created}, not an assumption about the
world. The $\E[X] = H_n$ bound holds no matter how the adversary arranges the
input, because we shuffle away their arrangement. This is the conversion promised
at the start of the chapter: assumption-about-input becomes
fact-enforced-by-randomness.

\begin{aestheticbox}
This pattern --- shuffle the input, then run a simple deterministic procedure ---
is everywhere. It is exactly how randomized quicksort defeats the
already-sorted-array worst case (Chapter on quicksort), how randomized hashing
defeats an adversary who tries to flood one bucket, and how load balancers spread
work. The lesson generalizes far past hiring: \emph{if a deterministic algorithm
has a rare bad case that an adversary could trigger, inject your own randomness
so that no input is reliably bad.}
\end{aestheticbox}

\section{Randomly Permuting an Array}
\label{ch05:permute}

That last algorithm leaned on a step we waved at: ``randomly permute the list.''
This is its own little algorithm and it is shockingly easy to get subtly wrong.
Our goal is a \emph{uniform random permutation}: a procedure that produces each
of the $n!$ orderings with probability exactly $1/n!$.

\subsection*{The right way: \textsc{Randomize-In-Place}}

The clean method shuffles the array in a single left-to-right pass. At step $i$
it swaps element $i$ with a randomly chosen element from positions $i$ through
$n$ (so possibly with itself).

\begin{algorithm}[H]
\caption{\textsc{Randomize-In-Place}$(A, n)$}
\KwIn{an array $A[1\dots n]$}
\KwOut{$A$ holds a uniform random permutation of its original contents}
\For{$i \gets 1$ \KwTo $n$}{
  $j \gets \textsc{Random}(i, n)$ \tcp*{uniform integer in $\{i,\dots,n\}$}
  swap $A[i] \leftrightarrow A[j]$\;
}
\end{algorithm}

Here \textsc{Random}$(a,b)$ returns an integer in $\{a, a+1, \dots, b\}$, each
equally likely, in $O(1)$ time. The whole routine is one loop of $n$ iterations
doing $O(1)$ work each, so it runs in $\Theta(n)$ time --- you cannot do better,
since you must touch every element.

The key picture is the \emph{shrinking window}: once position $i$ is filled it is
frozen, and the random draw for the next position is taken only from the unfrozen
tail. Figure~\ref{ch05:fig-shuffle} shows the window narrowing by one cell each
step for $n = 5$.

\begin{figure}[H]
\centering
\begin{tikzpicture}[
    cell/.style={draw,minimum width=0.85cm,minimum height=0.7cm},
    frozen/.style={cell,fill=targetgreen!18},
    window/.style={cell,fill=sourceblue!12},
  ]
  \foreach \row/\fcount/\lab in {0/0/{$i=1$}, 1/1/{$i=2$}, 2/2/{$i=3$}, 3/3/{$i=4$}, 4/4/{$i=5$}}{
    \node[left] at (-0.3,-\row*0.95) {\small \lab};
    \foreach \c in {1,2,3,4,5}{
      \pgfmathtruncatemacro{\isfrozen}{\c <= \fcount ? 1 : 0}
      \ifnum\isfrozen=1
        \node[frozen] at (\c*0.95,-\row*0.95) {\scriptsize fixed};
      \else
        \node[window] at (\c*0.95,-\row*0.95) {};
      \fi
    }
    % brace under the window (positions fcount+1 .. 5)
    \pgfmathsetmacro{\lo}{(\fcount+1)*0.95-0.45}
    \draw[decorate,decoration={brace,amplitude=3pt,mirror},sourceblue]
        (\lo,-\row*0.95-0.4) -- (5*0.95+0.45,-\row*0.95-0.4);
  }
  \node[sourceblue,right] at (5*0.95+0.7,-2*0.95-0.4)
     {\small draw $j\in\{i,\dots,n\}$ here};
  \draw[->,thick,quotientorange] (5*0.95+0.6,0.2) -- (5*0.95+0.6,-3.6)
     node[midway,right,quotientorange,align=left] {\small window\\\small shrinks};
\end{tikzpicture}
\caption{The shrinking pick-window of \textsc{Randomize-In-Place} for $n=5$.
Green cells are frozen (already chosen and never touched again); the blue tail is
the range $\{i,\dots,n\}$ from which \textsc{Random} draws at step $i$. Each step
freezes one more cell, so the window loses exactly one position per row.}
\label{ch05:fig-shuffle}
\end{figure}

This is sometimes called the Fisher--Yates shuffle, or Knuth's shuffle. Let us
see why it is uniform.

\begin{lemma}[Uniformity of \textsc{Randomize-In-Place}]
\label{ch05:lem-shuffle}
\textsc{Randomize-In-Place} produces a uniform random permutation: every ordering
appears with probability $1/n!$.
\end{lemma}

\begin{proof}[Proof by loop invariant]
We use the invariant: \emph{just before the iteration for index $i$, the subarray
$A[1\dots i-1]$ contains any given sequence of $i-1$ of the original elements with
probability $(n-i+1)!/n!$.} (Equivalently: the prefix of length $i-1$ is a
uniformly random arrangement of some $(i-1)$-subset.)

\textbf{Initialization ($i=1$).} The prefix $A[1\dots 0]$ is empty, so there is
exactly one ``arrangement'' of $0$ elements, and the claimed probability is
$(n-0)!/n! = 1$. True.

\textbf{Maintenance.} Assume the invariant holds before iteration $i$. Consider
any specific desired arrangement of $A[1\dots i]$, say $(x_1, x_2, \dots, x_i)$.
For $A[1\dots i]$ to end up as this sequence, two things must happen: the prefix
$A[1\dots i-1]$ already equaled $(x_1, \dots, x_{i-1})$ \emph{before} this
iteration (probability $(n-i+1)!/n!$ by the invariant), \emph{and} the swap put
$x_i$ into position $i$. Once the prefix is fixed, $x_i$ still sits somewhere in
$A[i \dots n]$ --- that is $n-i+1$ positions --- and \textsc{Random}$(i,n)$ picks
its position with probability $1/(n-i+1)$. These two are independent, so
\[
  \Pr\{A[1\dots i] = (x_1,\dots,x_i)\}
  = \frac{(n-i+1)!}{n!}\cdot\frac{1}{n-i+1}
  = \frac{(n-i)!}{n!},
\]
which is exactly the invariant with $i$ in place of $i-1$. Maintenance holds.

\textbf{Termination ($i = n+1$).} The invariant now describes the full array:
each complete permutation $(x_1, \dots, x_n)$ occurs with probability
$(n-n)!/n! = 1/n!$. That is uniformity.
\end{proof}

\begin{example}
\label{ch05:ex-shuffle-trace}
Trace the shuffle on $A = [\,a,\,b,\,c\,]$, $n=3$. Suppose \textsc{Random}
returns: $j=3$ at $i=1$, then $j=2$ at $i=2$, then $j=3$ at $i=3$.

\smallskip
\begin{tabular}{llll}
\toprule
$i$ & range for $j$ & drawn $j$ & array after swap \\
\midrule
start & --- & --- & $[a,b,c]$ \\
$1$ & $\{1,2,3\}$ & $3$ & $[c,b,a]$ \quad (swap $A[1]\leftrightarrow A[3]$) \\
$2$ & $\{2,3\}$ & $2$ & $[c,b,a]$ \quad (swap with self) \\
$3$ & $\{3\}$ & $3$ & $[c,b,a]$ \quad (forced) \\
\bottomrule
\end{tabular}

\smallskip
This run yields $[c,b,a]$. There are $3\cdot 2\cdot 1 = 6$ equally likely
sequences of draws ($3$ choices at $i=1$, then $2$, then $1$), and you can check
that the six runs land on the six distinct permutations exactly once each ---
which is precisely uniformity.
\end{example}

\subsection*{A tempting wrong way}

\begin{trapbox}
A natural-looking ``fix'' is to let $j$ range over \emph{all} of $\{1,\dots,n\}$
at every step instead of $\{i,\dots,n\}$:
\begin{quote}
\texttt{for $i = 1$ to $n$: swap $A[i]$ with $A[\textsc{Random}(1,n)]$}
\end{quote}
This looks more random --- every position can move anywhere! --- but it is
\emph{not} uniform. It performs $n$ swaps, each with $n$ equally likely choices,
so it follows one of $n^n$ equally likely paths. But there are only $n!$
permutations, and $n!$ does not divide $n^n$ for $n \ge 3$. By pigeonhole, the
permutations cannot all get equal probability: some orderings are slightly
favored. For $n=3$ the six permutations receive the path-counts
$\{4,4,4,5,5,5\}$ out of the $27$ paths --- not all $27/6 = 4.5$ --- so three
arrangements are over-represented and three are under-represented. The lesson: more
visible randomness is not the same as uniform randomness. Use the proven
\textsc{Randomize-In-Place} with its shrinking range.
\end{trapbox}

\begin{walkbox}
Take the broken shuffle above for $n=3$ on a walk and enumerate it. There are
$27$ equally likely $(j_1,j_2,j_3)$ triples. Tally which of the $6$ permutations
each triple produces. You should find three permutations appearing $5$ times and
three appearing $4$ times (so $3\cdot 5 + 3\cdot 4 = 27$ paths in all, as it must).
Then ask yourself: \emph{why} does the correct version,
with its shrinking range $\{i,\dots,n\}$, give exactly $3! = 6$ paths mapping
one-to-one onto the $6$ permutations? The contrast between ``$n!$ equally likely paths''
versus ``$n^n$ equally likely paths'' explains why the bias arises.
\end{walkbox}

\section{A Recursion-Tree-Free Look at the Numbers}

It helps to \emph{see} the harmonic sum that controls the hiring problem.
Each term $1/i$ is the probability that candidate $i$ breaks the record. The
early candidates almost always break it ($1/1, 1/2, 1/3, \dots$); the later ones
rarely do. Plotting the terms makes the $\ln n$ growth visible: it is the area
under the curve $1/x$.

\begin{center}
\begin{tikzpicture}[scale=1.0]
  % axes
  \draw[->,thick] (0,0) -- (7.3,0) node[right] {\small $i$};
  \draw[->,thick] (0,0) -- (0,3.2) node[above] {\small $\Pr\{\text{hire at } i\}=1/i$};
  % the 1/x curve
  \draw[thick,sourceblue,domain=0.34:7,smooth,samples=120] plot (\x,{2.6/\x});
  \node[sourceblue,right] at (4.5,1.1) {\small $y = 1/x$};
  % bars for 1/i at i = 1..6 (scaled by same 2.6)
  \foreach \i in {1,2,3,4,5,6}{
    \pgfmathsetmacro{\h}{2.6/\i}
    \draw[fill=green!18,draw=green!45!black] (\i-0.36,0) rectangle (\i+0.0,\h);
  }
  \foreach \i in {1,2,3,4,5,6}{ \node[below] at (\i-0.18,0) {\scriptsize \i}; }
  % brace + label
  \draw[decorate,decoration={brace,amplitude=5pt},orange!70!black]
        (0.2,-0.7) -- (6.2,-0.7);
  \node[orange!70!black,below] at (3.2,-0.85)
        {\small total area $\approx \sum 1/i = H_n \approx \ln n$};
\end{tikzpicture}
\end{center}

The bars are the actual probabilities $1/i$; the smooth curve is $1/x$. The sum
of the bar heights is $H_n$, and it tracks the area under $1/x$ from $1$ to $n$,
which is $\ln n$. That single picture is why ``a million applicants'' costs only
``fourteen hires.''

\section{A Second Helping: Expected Number of Inversions}
\label{ch05:inversions}

To prove the indicator method is a reusable hammer and not a one-off trick, let
us swing it at a different problem. An \emph{inversion} in an array $A[1\dots n]$
of distinct numbers is a pair of indices $(i,j)$ with $i < j$ but $A[i] > A[j]$
--- a pair that is out of order. Inversions measure how unsorted an array is
(insertion sort, from an earlier chapter, runs in time proportional to the number
of inversions). \emph{If the array is a uniform random permutation, how many
inversions do we expect?}

Before the algebra, a picture of what we are counting. Figure~\ref{ch05:fig-inversions}
shows a small array; every line connects a pair $(i,j)$ with $i<j$, and a pair is
an inversion exactly when the left value is the larger one.

\begin{figure}[H]
\centering
\begin{tikzpicture}[
    val/.style={draw,circle,minimum size=0.7cm,inner sep=0pt},
  ]
  % array A = [3,1,4,2] at positions 1..4
  \foreach \pos/\v in {1/3,2/1,3/4,4/2}{
    \node[val] (p\pos) at (\pos*1.6,0) {\small $\v$};
    \node[below=2pt of p\pos] {\scriptsize pos \pos};
  }
  % inversion arcs (i<j, A[i]>A[j]): (1,2):3>1, (1,4):3>2, (3,4):4>2
  \draw[trapred,thick,bend left=45] (p1) to (p2);
  \draw[trapred,thick,bend left=45] (p1) to (p4);
  \draw[trapred,thick,bend left=40] (p3) to (p4);
  % non-inversion arcs (in order): (1,3):3<4, (2,3):1<4, (2,4):1<2
  \draw[targetgreen,thick,dashed,bend right=35] (p1) to (p3);
  \draw[targetgreen,thick,dashed,bend right=35] (p2) to (p3);
  \draw[targetgreen,thick,dashed,bend right=40] (p2) to (p4);
\end{tikzpicture}
\caption{Inversions of $A=[3,1,4,2]$. Solid red arcs above are the three
inversions ($3{>}1$, $3{>}2$, $4{>}2$); dashed green arcs below are the three
in-order pairs. There are $\binom{4}{2}=6$ pairs total, and here exactly half are
inverted --- the average outcome the calculation below predicts.}
\label{ch05:fig-inversions}
\end{figure}

Run the recipe. For each pair $i < j$ define the indicator
\[
  X_{ij} = I\{\,A[i] > A[j]\,\} = I\{(i,j)\text{ is an inversion}\}.
\]
The total number of inversions is $X = \sum_{i<j} X_{ij}$. By linearity,
\[
  \E[X] = \sum_{1 \le i < j \le n} \E[X_{ij}]
        = \sum_{1 \le i < j \le n} \Pr\{A[i] > A[j]\}.
\]
For any fixed pair of distinct values, in a uniform random permutation they are
equally likely to appear in either order, so $\Pr\{A[i] > A[j]\} = \tfrac12$. The
number of pairs with $i<j$ is $\binom{n}{2} = \tfrac{n(n-1)}{2}$. Therefore
\[
  \E[X] = \binom{n}{2}\cdot\frac12 = \frac{n(n-1)}{4}.
\]

\begin{example}
For $n = 10$: $\E[X] = \tfrac{10\cdot 9}{4} = 22.5$ inversions on average. For
$n = 100$: $\tfrac{100\cdot 99}{4} = 2475$. Roughly a quarter of all pairs are
inverted in a random array --- which is why insertion sort, costing one unit per
inversion, runs in expected $\Theta(n^2)$ time on random input. The indicator
method just \emph{told} us the average-case running time of an algorithm from
two earlier chapters, almost for free.
\end{example}

\begin{focusbox}
Look back at the three problems --- heads in coin flips, hires, inversions. The
\emph{same four-line recipe} solved all three. Decompose into indicators, apply
linearity, swap each expectation for a probability, sum. When you see ``expected
number of \emph{[somethings]},'' your hand should already be reaching for an
indicator per \emph{something}.
\end{focusbox}

\section{Las Vegas vs.\ Monte Carlo: Two Flavors of Randomized Algorithm}

Randomized algorithms come in two temperaments, and it is worth naming them.

\begin{intuitionbox}
A \textbf{Las Vegas} algorithm always returns the \emph{correct} answer; only its
\emph{running time} is random. Randomized quicksort and \textsc{Randomized-Hire}
are Las Vegas: the output is never wrong, but how long it takes depends on the
coins. You gamble with time, never with correctness.

A \textbf{Monte Carlo} algorithm always finishes in a \emph{bounded} time, but
may occasionally return a \emph{wrong} answer, with some controllable small
probability. Many primality tests and approximate-counting schemes are Monte
Carlo: fast and bounded, but right only ``with high probability.''

Mnemonic: in Las Vegas you are guaranteed to leave with the right answer, you
just do not know \emph{when}; in Monte Carlo you leave \emph{on time}, but might
leave with the wrong souvenir.
\end{intuitionbox}

This chapter's algorithms are all Las Vegas. We will meet Monte Carlo algorithms
later, especially in the number-theory chapter, where being wrong one time in a
billion is a perfectly good trade for enormous speed.

\section{What ``Expected Running Time'' Really Means}

We should be careful about a phrase we have been using loosely. When we say a
randomized algorithm has \emph{expected running time} $T(n)$, we mean: \emph{fix}
any input of size $n$, then average the running time over the algorithm's
internal coin flips. The expectation is taken over the randomness of the
algorithm, with the input held fixed and adversarial.

\begin{formalbox}
For a randomized algorithm $\mathcal{A}$, the expected running time on an input
$x$ is
\[
  T_{\mathcal{A}}(x) = \E_{r}\big[\, \text{steps } \mathcal{A} \text{ takes on
  input } x \text{ using random bits } r \,\big],
\]
where the expectation is over the algorithm's random bits $r$. The
\emph{worst-case expected running time} is then $\max_{|x| = n}
T_{\mathcal{A}}(x)$: the worst this average gets over all inputs of size $n$.
\end{formalbox}

\begin{trapbox}
Do not confuse \emph{expected} running time with \emph{worst-case} running time.
Randomized quicksort has expected running time $\Theta(n \log n)$ but worst-case
running time $\Theta(n^2)$ --- the bad case still exists, it has just become
astronomically unlikely (the coins would all have to conspire). ``Expected
$\Theta(n\log n)$'' is a promise about the average over coins, not a promise that
$n^2$ can never happen.
\end{trapbox}

Figure~\ref{ch05:fig-pivot-tree} shows this intuition through the probability tree of a
randomized pivot choice. When a pivot is chosen uniformly at random, it lands in the
central half of the sorted array with probability $1/2$. Whenever it lands there (a ``good split''),
neither recursive subarray exceeds $\frac{3}{4}n$ in size. Even if a pivot lands in the outer
quarters (a ``bad split''), the next random trial is an independent coin flip with success
probability $1/2$. The expected number of trials until a good split is at most $1/(1/2) = 2$.
On average, every two partition steps shrink the problem by a factor of at least $3/4$,
capping the expected recursion depth at $O(\log n)$.

\begin{figure}[H]
\centering
\begin{tikzpicture}[
    level distance=14mm,
    level 1/.style={sibling distance=38mm},
    level 2/.style={sibling distance=26mm},
    every node/.style={font=\small},
    pbox/.style={draw,rounded corners,inner sep=3.5pt},
    prob/.style={font=\scriptsize,sourceblue!80!black}
  ]
  % Partition bar above tree
  \begin{scope}[yshift=2.0cm, xshift=-1.6cm]
    \node[anchor=east,font=\footnotesize] at (-0.3, 0.25) {Pivot ranks:};
    \draw[fill=trapred!15] (0,0) rectangle (1.2, 0.5) node[midway,font=\tiny] {$\le n/4$};
    \draw[fill=targetgreen!20,thick] (1.2,0) rectangle (3.6, 0.5) node[midway,font=\footnotesize\bfseries] {Good: $n/4 \dots 3n/4$};
    \draw[fill=trapred!15] (3.6,0) rectangle (4.8, 0.5) node[midway,font=\tiny] {$> 3n/4$};
    \node[font=\scriptsize,targetgreen!50!black] at (2.4, 0.8) {probability $= 1/2$};
  \end{scope}

  % Tree
  \node[pbox,fill=sourceblue!10] (root) {Input size $n$}
    child {
      node[pbox,fill=targetgreen!18] {Good: $\le \tfrac{3}{4}n$}
      edge from parent node[left,prob] {$p = \tfrac{1}{2}$}
    }
    child {
      node[pbox,fill=trapred!14] (bad1) {Bad: $\approx n$}
      child {
        node[pbox,fill=targetgreen!18] {Good: $\le \tfrac{3}{4}n$}
        edge from parent node[left,prob] {$p = \tfrac{1}{2}$}
      }
      child {
        node[pbox,fill=trapred!14] {Bad: $\approx n$}
        edge from parent node[right,prob] {$p = \tfrac{1}{2}$}
      }
      edge from parent node[right,prob] {$p = \tfrac{1}{2}$}
    };
\end{tikzpicture}
\caption{Randomized pivot partitioning represented as a probability tree. Half of all
possible pivot choices fall in the central $50\%$ of the ranks, guaranteeing that the larger
subproblem has size at most $\frac{3}{4}n$. A geometric waiting time shows that a good split
occurs within $2$ steps in expectation, yielding $O(\log n)$ expected tree depth.}
\label{ch05:fig-pivot-tree}
\end{figure}

\section{Putting It Together: a Tiny End-to-End Example}

Let us run one complete numeric story through the machinery, hiring problem
style, to leave no gap between the abstract recipe and concrete arithmetic.

\begin{example}[Hiring with $n = 4$, fully enumerated]
\label{ch05:ex-n4}
Take four candidates with qualities $1 < 2 < 3 < 4$ (so $4$ is best). With a
uniform random arrival order there are $4! = 24$ equally likely orderings. We
claim $\E[\text{hires}] = H_4 = 1 + \tfrac12 + \tfrac13 + \tfrac14 =
\tfrac{25}{12} \approx 2.083$.

Let us sanity-check by counting hires for a few orders (candidate $0$ is the
dummy, so the first real candidate is always hired):
\begin{itemize}
  \item Order $1,2,3,4$: hire at every step. $4$ hires.
  \item Order $4,3,2,1$: hire only the first. $1$ hire.
  \item Order $2,1,4,3$: hire $2$, then $4$. $2$ hires.
  \item Order $3,1,4,2$: hire $3$, then $4$. $2$ hires.
\end{itemize}
If you grind through all $24$ orders, the hire counts sum to $50$, giving average
$50/24 = 25/12 \approx 2.083$ --- exactly $H_4$. The indicator argument gave this
in four lines; the brute-force enumeration confirms it. Seeing both agree is the
point: the recipe is not magic, it is bookkeeping done cleverly.
\end{example}

\begin{labbox}
Reproduce Example~\ref{ch05:ex-n4} yourself. Write out all $24$ permutations of
$\{1,2,3,4\}$ (or just generate them mentally in lexicographic order), count
hires for each, and confirm the total is $50$. Then group the $24$ orders by
``did candidate at position $i$ get hired?'' and verify that position $i$ is a
hire in exactly $24/i$ of the orders --- $24$ for $i=1$, $12$ for $i=2$, $8$ for
$i=3$, $6$ for $i=4$ --- which is the indicator probability $1/i$ in disguise.
\end{labbox}

\begin{complexitybox}
\textbf{Costs at a glance.}
\begin{itemize}
  \item \textsc{Randomize-In-Place}: $\Theta(n)$ time, $O(1)$ extra space;
        produces each of $n!$ permutations with probability $1/n!$.
  \item \textsc{Hire-Assistant} (deterministic, random input): interviewing
        $\Theta(n)$ always; expected hires $H_n = \ln n + O(1)$; worst-case hires
        $n$.
  \item \textsc{Randomized-Hire-Assistant}: same expected hires $H_n$ for
        \emph{every} input (adversary cannot force the worst case); cost of the
        shuffle is $\Theta(n)$, dwarfed by interviewing.
  \item Expected inversions in a uniform random permutation:
        $\binom{n}{2}/2 = n(n-1)/4 = \Theta(n^2)$.
\end{itemize}
\end{complexitybox}

\begin{historybox}
In 1946, while recovering from encephalitis and playing Canfield solitaire, Polish-American
mathematician Stanislaw Ulam wondered what fraction of solitaire deals are theoretically winnable.
Working out the answer by pure combinatorics was an impenetrable nightmare. Ulam had an epiphany:
instead of solving the unsolvable equation, why not simply deal one hundred random hands, count the
wins, and observe the empirical frequency? He shared the idea with John von Neumann at Los Alamos,
who immediately saw its life-or-death value for simulating neutron transport and nuclear fission in
the Manhattan Project. In 1947, von Neumann implemented the first large-scale automated simulation
on the vacuum-tube ENIAC, codenaming it ``Monte Carlo'' after the Monaco gambling casino. Three
decades later, in 1976, Michael Rabin published his foundational paper ``Probabilistic Algorithms,''
demonstrating that tossing coins inside an algorithm could tame notoriously hard discrete problems
--- from testing whether a thousand-digit number is prime to finding the closest pair of points in
space.

The rise of randomized algorithms met fierce philosophical resistance. For generations, scientists
believed an algorithm that could produce different running times or answers on the same input was
broken or untrustworthy. It took decades to distinguish between \emph{Monte Carlo} algorithms (which
run in fixed time but carry an infinitesimal, tunable chance of error) and \emph{Las Vegas}
algorithms (which never return an incorrect answer, but whose running time is a random variable).
Early implementations were also plagued by bad pseudo-random generators; von Neumann's middle-square
method frequently decayed into zero, while naive array shuffles (swapping each $A[i]$ with
$A[\textsc{Random}(1, n)]$) explored $n^n$ paths, which cannot divide $n!$ evenly and thus quietly
rigged card games in early software.

Randomization unlocked an unbeatable defense against malicious adversaries: when an algorithm flips
its own coins at runtime, no single worst-case input can ever be crafted in advance. It powers
universal hashing to thwart denial-of-service attacks and enables polynomial-time prime generation
for RSA encryption. Yet its honest downsides are real: software requires access to genuine physical
entropy pools (such as thermal jitter), and tracking down a bug that only triggers once every ten
million random coin flips can confound even the best engineers. In Chapter~7, we apply these coin
flips directly to pivot selection in Randomized Quicksort; in Chapter~11, we use random hash families
to build collision-resistant hash tables; and in Chapter~30, we see Rabin's primality test in full
mathematical glory.
\end{historybox}

\begin{aestheticbox}
Why does any of this matter outside a textbook? Because the ``shuffle, then run
something simple'' pattern is load-bearing in real systems. Randomized quicksort
is the default sort in many language standard libraries precisely because
randomization makes the $\Theta(n^2)$ case vanish in practice without any input
assumptions. Hash tables randomize their hash function so that no adversary can
craft keys that all collide --- a real security concern (``hash-flooding''
denial-of-service attacks against web servers were patched exactly this way).
Distributed systems pick random servers to balance load. Cryptography is
\emph{built} on randomness. The humble indicator random variable is your first
rigorous tool for predicting how all of these behave on average.
\end{aestheticbox}

\section{Summary}

You arrived with deterministic machines and leave with dice in your toolbox.
The throughline:

\begin{itemize}
  \item \textbf{Two averages.} Probabilistic analysis averages a deterministic
        algorithm over a random \emph{input}; a randomized algorithm averages
        over its own \emph{coin flips}. Prefer the second: it guarantees the
        bound for every input.
  \item \textbf{Expectation and its linearity.} $\E[X] = \sum_x x\Pr\{X=x\}$,
        and $\E[\sum X_i] = \sum \E[X_i]$ \emph{always}, correlated or not. That
        no-fine-print linearity is the workhorse.
  \item \textbf{Indicator random variables.} $\E[I\{A\}] = \Pr\{A\}$. To find an
        expected count, write it as a sum of indicators, apply linearity, replace
        each expectation by a probability, and sum. We ran this recipe on coin
        flips ($n/2$), hires ($H_n$), and inversions ($n(n-1)/4$).
  \item \textbf{The hiring problem.} Expected hires $= H_n = \ln n + O(1)$:
        tiny. Randomizing the input order makes this guarantee hold against any
        adversary.
  \item \textbf{Uniform shuffling.} \textsc{Randomize-In-Place} swaps $A[i]$ with
        a random later (or equal) position, runs in $\Theta(n)$, and is provably
        uniform by a loop invariant. The all-positions variant looks more random
        but is provably \emph{not} uniform.
\end{itemize}

Everything from here --- quicksort, hashing, skip lists, primality testing ---
leans on these few ideas. Hold onto the indicator recipe most tightly; it is the
single most reused tool in the analysis of algorithms.

\section{Exercises}

\begin{exercise}[Warm-up: expectation of a biased coin]
A coin comes up heads with probability $p$ and tails with probability $1-p$. Let
$X$ be the indicator that a single flip is heads. Compute $\E[X]$ directly from
the definition of expectation, and confirm it equals $\Pr\{\text{heads}\}$ as
Lemma~\ref{ch05:lem-indicator} promises. Then, for $n$ independent such flips,
find the expected number of heads using indicators.
\end{exercise}

\begin{exercise}[Expected maximum index of a hire]
In \textsc{Hire-Assistant} with random input, candidate $1$ is always hired and
candidate $n$ is hired with probability $1/n$. Using indicators, what is the
expected number of \emph{interviews that result in a hire among the last half} of
candidates, i.e.\ positions $\lceil n/2\rceil$ through $n$? Express your answer
as a difference of harmonic numbers and estimate it for large $n$.
\end{exercise}

\begin{exercise}[Two dice]
Roll two fair dice and let $X$ be the sum of the faces. Use linearity of
expectation (write $X = X_1 + X_2$ for the two dice) to find $\E[X]$ in two
lines, without summing over all $36$ outcomes. Then find $\E[X]$ the long way and
confirm they agree.
\end{exercise}

\begin{exercise}[Tracing the shuffle]
Run \textsc{Randomize-In-Place} on $A = [w,x,y,z]$ ($n=4$) assuming
\textsc{Random} returns $j = 2, 4, 3, 4$ at steps $i = 1,2,3,4$ respectively.
Show the array after each swap. Then argue, using
Lemma~\ref{ch05:lem-shuffle}, why this particular output had probability exactly
$1/24$ before the run started.
\end{exercise}

\begin{exercise}[Why the broken shuffle fails, in miniature]
Consider the incorrect shuffle that, for each $i$ from $1$ to $n$, swaps $A[i]$
with $A[\textsc{Random}(1,n)]$. For $n = 2$, enumerate all $2^2 = 4$ equally
likely paths and compute the probability of each of the two permutations of
$[a,b]$. Are they both $1/2$? Now do the same for the \emph{correct}
\textsc{Randomize-In-Place} and confirm it gives $1/2$ each. (For $n=2$ the
broken version happens to be uniform; explain why it first breaks at $n=3$, in
terms of $n!$ dividing $n^n$.)
\end{exercise}

\begin{exercise}[Expected inversions, restricted]
In a uniform random permutation of $\{1,\dots,n\}$, what is the expected number
of inversions $(i,j)$ with $i<j$ that are \emph{adjacent}, i.e.\ $j = i+1$? Use
one indicator per adjacent pair, and compare your answer to the total expected
inversions $n(n-1)/4$. What fraction of inversions, on average, are adjacent?
\end{exercise}

\begin{exercise}[Expected number of fixed points]
A \emph{fixed point} of a permutation $\pi$ of $\{1,\dots,n\}$ is an index $i$
with $\pi(i) = i$. For a uniform random permutation, define $X_i = I\{\pi(i) =
i\}$ and use the indicator recipe to show the expected number of fixed points is
exactly $1$, \emph{independent of $n$}. (This is a famous and counterintuitive
fact: whether you shuffle $5$ cards or $5$ million, on average exactly one ends
up in its original place.)
\end{exercise}

\begin{exercise}[Coupon collector, gently]
You repeatedly draw a uniformly random coupon from $n$ distinct types (with
replacement). Let $T$ be the number of draws until you have collected all $n$
types. Split the process into $n$ \emph{phases}: phase $k$ is the time spent
holding exactly $k-1$ distinct types while waiting for a new one. Argue that the
expected length of phase $k$ is $\frac{n}{n-k+1}$ (a geometric waiting time), then
use linearity to show $\E[T] = n H_n \approx n\ln n$. Compare this $n\ln n$ to the
hiring problem's $\ln n$ and explain in words why one carries an extra factor of
$n$.
\end{exercise}

\begin{exercise}[Challenge: expected comparisons in a max-scan]
You scan an array of $n$ distinct numbers in random order, keeping a running
maximum (like \textsc{Hire-Assistant}, but you only update a variable rather than
pay a hiring fee). Each time you see a new record you do one ``update.'' You
already know the expected number of updates is $H_n$. Now suppose each update
costs $i$ units when it happens at position $i$ (later updates are more
expensive). Using indicators \emph{and} the fact that $\E[i \cdot X_i] = i \cdot
\E[X_i]$ here (why is this valid even though $i$ multiplies the indicator?),
compute the expected total update cost $\sum_{i=1}^n i \cdot \Pr\{\text{record at
} i\}$. Simplify it to a closed form and give its asymptotic order.
\end{exercise}

\begin{exercise}[Challenge: Monte Carlo majority]
An array of $n$ elements is promised to contain a \emph{strict majority} element
(one value occupying more than $n/2$ slots). Here is a Monte Carlo algorithm:
pick a random index, read its value $v$, and \emph{guess} that $v$ is the
majority. (a) What is the probability this single guess is correct? (b) If you
repeat the guess $k$ times independently and output any value that you guessed at
least once, bound the probability that the true majority is never guessed.
(c) How large must $k$ be to make the failure probability below $1/1000$? This is
the essence of probability amplification --- turning a coin-flip-quality
algorithm into a reliable one by repetition.
\end{exercise}
